% ============================================================
%  Chapter 2: Background
%  Meant to be included from main.tex via \input{background}.
%  Do NOT add \documentclass or \begin{document} here.
%  Assumes IEEEtran bibliography style and biblatex-or-bibtex
%  configured in main.tex.
% ============================================================

\chapter{Background}
\label{chap:background}

This chapter defines the structured-extraction concepts the rest of
the dissertation relies on; critical comparison of competing
approaches is deferred to Chapter~\ref{chap:litreview} and the
evaluation metrics used to score extractions are defined in the
Evaluation chapter where they are applied.

% ------------------------------------------------------------
\section{Document intelligence and structured data extraction}
\label{sec:bg-docintel}

Document intelligence converts human-readable documents into
machine-readable representations. A typical pipeline takes a rendered
page (a scan, a photograph, or a native PDF) and produces structured
records bound to a schema: named entities, typed fields, tables, or
key-value pairs~\cite{cui2021documentai}. The core difficulty is that
entity vocabularies and page structures vary across domains and
languages~\cite{silva2021review}.

Two benchmarks recur throughout this dissertation.
CORD~\cite{park2019cord} is a consolidated receipt dataset with
per-box text and multi-level semantic labels.
DocVQA~\cite{mathew2021docvqa} reframes document understanding as
visual question answering over document images. Both measure success
by how well the produced record matches a reference one, not by
full-page transcription fidelity.

% ------------------------------------------------------------
\section{Semi-structured documents and utility bills}
\label{sec:bg-semistructured}

Documents are informally divided into three classes. Unstructured
documents (narrative prose, articles) have no reliable positional
convention. Structured documents (relational database exports, filled
digital forms with named fields) carry their schema explicitly.
Semi-structured documents sit between the two: the same set of fields
recurs across instances, but layout, font, language, and ordering
change from one issuer to the next~\cite{abiteboul1997querying}.

\begin{figure}[ht]
\centering
\includegraphics[width=0.95\linewidth]{Chapters/figures/fig-2-1-document-class-spectrum.png}
\caption[Document class spectrum.]{The three-class document spectrum.
Utility bills sit in the semi-structured middle: the schema is
stable, the surface layout varies by issuer.}
\label{fig:bg-doc-spectrum}
\end{figure}

Receipts and invoices are the canonical examples.
CloudScan~\cite{palm2017cloudscan} established cross-issuer
generalisation as a distinct evaluation axis for invoice field
extraction, separating performance on previously-seen layouts from
performance on unseen ones. ReceiptSense~\cite{abdallah2024receiptsense}
shows that the same class of document exhibits multilingual surface
variation at scale, motivating the multilingual design of the corpus
used here.

Utility bills belong to the same family. Every issuer prints broadly
the same information (account holder, billing period, consumption,
amount due, due date) because these fields are mandated by sector
regulation; Article~18 of the EU electricity market directive binds
customer bills to the minimum content requirements set out in Annex~I
of that directive~\cite{eu2019electricity}. Despite this common core,
there is no cross-issuer standard for where on the page any of it
appears, what label is used, what currency or language the bill is in,
or what typographic conventions govern totals and subtotals. The
schema is fixed and small enough to define crisply; the surface
variation is wide.

\begin{figure}[ht]
\centering
\includegraphics[width=0.95\linewidth]{Chapters/figures/fig-2-2-utility-bill-layouts.png}
\caption[Same schema, different surfaces.]{The same five schema
fields appear at different positions, under different labels, and
in different languages across three mock issuers. The schema is
fixed; the surface varies.}
\label{fig:bg-bill-layouts}
\end{figure}

% ------------------------------------------------------------
\section{The template-based extraction paradigm}
\label{sec:bg-template}

Before the current generation of learned models, structured extraction
from semi-structured documents was approached by binding rules to page
regions. Klink and Kieninger~\cite{klink2001rule} proposed a hybrid
system in which geometrical layout features and textual cues combine
through fuzzy-matched rules, with domain-specific sets layered on top
of domain-independent routines. The deployment pattern that grew out
of this work is what the document-analysis literature calls a
\emph{template}~\cite{schuster2013intellix,esser2014fewexemplar}: a
hand-written or semi-automatically induced rule set keyed to a
specific issuer's layout, to be re-authored whenever the layout
changes.

A second wave of template-style methods replaced explicit rules with
2D-grid neural representations. Chargrid~\cite{katti2018chargrid}
encodes a document as a 2D grid of one-hot character vectors aligned
with page coordinates and passes the grid through a fully convolutional
encoder-decoder to predict segmentation masks and bounding boxes for
target fields. BERTgrid~\cite{denk2019bertgrid} extends this by
replacing the character encoding with contextualised BERT word-piece
embeddings, so the spatial representation carries language semantics.
Both are trained per field schema. Majumder et
al.~\cite{majumder2020form} extend the family with representation
learning over extraction candidates rather than whole-page layout.

% ------------------------------------------------------------
\section{Optical character recognition}
\label{sec:bg-ocr}

Optical character recognition (OCR) converts an image of text into a
machine-readable character sequence, typically with bounding boxes.
Modern engines combine a text-detection stage (where on the page is
there text) with a text-recognition stage (what does it say), an
architecture Long et al.~\cite{long2021scenetext} identify as dominant
across the deep-learning era. Tesseract~\cite{smith2007tesseract} is
the reference open-source engine; its 2007 overview describes line
finding, feature-based classification, and an adaptive classifier. It
is the default OCR stage in academic pipelines that need local
execution, multilingual language-pack support, and zero per-document
cost.

Two properties of OCR matter for downstream extraction. Errors are
not uniformly distributed: they concentrate on historical and degraded
print and on multilingual text whose vocabulary drifts from the
dictionaries used for correction, as catalogued in the post-OCR
processing survey of Nguyen et al.~\cite{nguyen2021postocr}.
Lopresti~\cite{lopresti2009ocrerrors} traces how OCR noise cascades
through sentence-boundary detection, tokenisation, and part-of-speech
tagging: an early error can disable every subsequent stage that
depends on clean tokens. Any pipeline that places an OCR stage before
a reasoning stage inherits the OCR stage's error profile.

\begin{figure}[ht]
\centering
\includegraphics[width=0.95\linewidth]{Chapters/figures/fig-2-3-ocr-error-cascade.png}
\caption[OCR error cascade.]{OCR errors propagate: a single
recognition mistake at stage 1 disables every subsequent stage that
depends on clean tokens.}
\label{fig:bg-ocr-cascade}
\end{figure}

% ------------------------------------------------------------
\section{Large language models for information extraction}
\label{sec:bg-llms}

A large language model (LLM) is an autoregressive neural network
trained to predict the next token on a large text
corpus~\cite{zhao2023llm}, built on the Transformer of Vaswani et
al.~\cite{vaswani2017attention} which replaced recurrence with
self-attention and enabled parallel training at previously
impractical scales. Brown et al.~\cite{brown2020gpt3} showed that at
sufficient scale these models perform many tasks from a few in-prompt
examples without gradient updates; Ouyang et
al.~\cite{ouyang2022instructgpt} then demonstrated that instruction
tuning with human feedback aligns model outputs with natural-language
task descriptions. The latter capability is what this dissertation
relies on when it specifies extraction behaviour through a prompt
rather than through training.

Applied to information extraction, an LLM reads an input (typically
text, possibly the output of an OCR engine) and emits a structured
record. The naive realisation is to prompt for JSON and parse the
result, which fails when the model emits syntactically invalid JSON,
a field name outside the schema, or a value of the wrong type. Two
mechanisms address the reliability problem. Constrained decoding
intervenes at generation time by masking any token that would violate
the target grammar; Willard and Louf~\cite{willard2023guided}
formalise this as finite-state machine transitions over the
vocabulary, with regular expressions and context-free grammars
compiled into the mask. Schema-enforced provider APIs expose an
equivalent mechanism at the service level. OpenAI's Structured
Outputs~\cite{openai2024structured} accept a JSON Schema and constrain
the decoder so that any returned string parses against the schema,
with 100\,\% schema adherence reported for
\texttt{gpt-4o-2024-08-06} under Structured Outputs, against under
40\,\% for \texttt{gpt-4-0613} without it.

One failure mode generalises across text- and vision-mode extraction.
\emph{Hallucination} is the production of a plausible-looking output
value that is not supported by the input; Ji et
al.~\cite{ji2023hallucination} survey it across generative NLP tasks
as a systematic rather than occasional failure. In document
extraction, hallucination appears as a value that fits the field's
type and looks appropriate for the document class but does not occur
in the source.

For the rest of this dissertation, an LLM extractor means a model
that reads tokens and produces a record conforming to a supplied
schema. Whether the input tokens come from an OCR engine or from
native text is the separate design decision treated in
Chapter~\ref{chap:litreview}.

% ------------------------------------------------------------
\section{Vision-language models and multimodal input}
\label{sec:bg-vlms}

A vision-language model (VLM) is an LLM extended to accept image
input alongside text. Most VLMs connect a vision encoder to a
language model through a learned projection and are then
instruction-tuned on multimodal data. Yin et al.~\cite{yin2024mllm}
survey this ``vision encoder + projection + LLM'' pattern as the
dominant design.

\begin{figure}[ht]
\centering
\includegraphics[width=0.9\linewidth]{Chapters/figures/fig-2-4-vlm-architecture.png}
\caption[VLM architecture.]{The dominant VLM pattern: a vision
encoder, a projection layer, and a language decoder, accepting
images alongside text.}
\label{fig:bg-vlm-arch}
\end{figure}

Two families of VLM are relevant. The first is the
\emph{general-purpose} frontier models, which accept arbitrary images
alongside text prompts, were not trained specifically for documents,
and are reached through provider APIs rather than self-hosted. The
two evaluated in this dissertation are OpenAI's
GPT-4o~\cite{openai2024gpt4o} and Anthropic's
Claude~\cite{anthropic2024claude}. The second is the
\emph{specialised} document models trained end-to-end on document
images. Donut~\cite{kim2022donut} is the prototypical example: it
eliminates the external OCR stage entirely, mapping a document image
to a structured output sequence via a transformer encoder-decoder
pretrained with a synthetic document generator for language and
domain coverage. LayoutLMv3~\cite{huang2022layoutlmv3} pretrains a
multimodal transformer with unified text and image masking plus a
word-patch alignment objective, reporting state-of-the-art results on
form understanding, receipt understanding, document VQA, document
classification, and layout analysis.

The VLM-specific form of the hallucination problem introduced in
\S\ref{sec:bg-llms} is known in the vision literature as \emph{object
hallucination}: the production of output describing entities not
present in the input image. Li et al.~\cite{li2023pope} find it
widespread across representative VLMs. Vision-mode extraction
therefore avoids the OCR error cascade of \S\ref{sec:bg-ocr} but
admits a visual variant of the hallucination risk defined in
\S\ref{sec:bg-llms}.

\begin{figure}[ht]
\centering
\includegraphics[width=\linewidth]{Chapters/figures/fig-2-5-two-path-pipeline.png}
\caption[Two extraction pipelines compared in this dissertation.]{The
two extraction pipelines this dissertation compares. Both share the
rasteriser, the LLM, and the output schema. The only variable is
how pixels become model input.}
\label{fig:bg-two-path}
\end{figure}
